\chapter{Read the actual model code, one step at a time}
\label{chap:code}

\section{The map before the details}
The file \code{src/MoranModel.jl} contains the rules and calculations. It does not run the whole study simply because you open it. A script in the \code{scripts} folder loads those rules and calls the functions. This chapter explains the model file. Chapter~10 explains the scripts that organise the experiments and save the results.

There are nine functions. First we reject impossible inputs. Next we calculate the probabilities for one event. Then we turn a random number into an event. Those three pieces let us build a complete history and repeat the experiment. The remaining functions calculate an exact answer, check that answer another way, calculate a confidence interval, and produce a deterministic comparison.

The listings below use byte-for-byte copies of the runnable files, refreshed by the guide-preparation script. Keeping these copies inside the guide folder allows the LaTeX editor to compile a self-contained project. The numbers at the left are its actual line numbers. Appendix~\ref{app:source} contains the complete file without interruptions.

\lstinputlisting[firstline=1,lastline=8,firstnumber=1]{code/MoranModel.jl}

Line~1 begins a module named \code{MoranModel}. Lines~3--4 load Julia's standard tools for random numbers and linear algebra. Linear algebra is a way of solving several linked equations; we explain our small example in Chapter~7. Lines~6--8 export the names of our functions, making them convenient to call after loading the module. The final \code{end} at the bottom of the file closes this module. It is separate from the \code{end} that closes each function.

\section{Function 1: reject inputs that do not describe our model}
\lstinputlisting[firstline=11,lastline=16,firstnumber=11]{code/MoranModel.jl}

\begin{description}[style=nextline]
\item[Line 11] Define a function receiving population size \code{N}, focal count \code{i}, and selective advantage \code{s}.
\item[Line 12] \code{isa Integer} checks for a whole-number input. The double ampersand means ``and.'' Both the integer check and \code{N >= 2} must be true. The double vertical bar means ``or'': if the preceding condition fails, run \code{throw}, which stops with a descriptive error. \code{ArgumentError} names the kind of error. The biological requirement is a population of at least two individuals.
\item[Line 13] Apply the same idea to the allele count. Zero is allowed because the allele may already be lost; \code{N} is allowed because it may already be fixed.
\item[Line 14] Require a finite real number at least zero. \code{isfinite} excludes infinity and undefined numerical values. A positive input must also remain finite and positive when converted to \code{Float64}; we reject numbers outside this implementation's representable range. Negative selection is outside this teaching project's scope.
\item[Line 15] \code{nothing} is Julia's value meaning that there is no ordinary result to return. This function checks inputs; it does not calculate a probability.
\item[Line 16] Close the function. A passed check allows the caller to continue; a failed check prevents it from silently producing an answer for the wrong model.
\end{description}

\section{Function 2: calculate the three event probabilities}
\lstinputlisting[firstline=19,lastline=32,firstnumber=19]{code/MoranModel.jl}

The mathematical derivation used total reproductive weight \((1+s)i+(N-i)\). The code uses a harmless rescaling. Dividing \emph{every} individual's reproductive weight by the same positive number does not change which type is more likely to reproduce. In particular,
\[
\frac{(1+s)i}{(1+s)i+(N-i)}
=\frac{i}{i+(N-i)/(1+s)}.
\]
We divided both the top and the bottom by \(1+s\). For \(N=5,i=2,s=0.5\), the original weights are \(3\) and \(3\); the scaled weights are \(2\) and \(2\). Either calculation gives a parent probability of one half.

\begin{description}[style=nextline]
\item[Lines 19--20] Begin the function and check its inputs.
\item[Line 21] Calculate \(1+s\). \code{Float64} converts the input to Julia's ordinary double-precision floating-point representation. The \code{1.0} is a floating-point one.
\item[Line 22] Explain the rescaling in a comment. It avoids multiplying a large fitness by the count.
\item[Line 23] Give each focal individual scaled weight one, so their combined weight is \code{i}.
\item[Line 24] Give each other individual scaled weight \(1/(1+s)\), so their combined weight is \((N-i)/(1+s)\).
\item[Line 25] Add the two weights to find the total.
\item[Lines 26--27] Divide each type's weight by the total. These are parent probabilities, not probabilities of changing the count.
\item[Line 28] To increase the count, a focal parent must reproduce \emph{and} an individual of the other type must be replaced. Multiply the first probability by \((N-i)/N\).
\item[Line 29] To decrease the count, the other type must reproduce and a focal individual must be replaced. Multiply the relevant probabilities.
\item[Line 30] Whatever probability remains is the chance of no change. The three outcomes cover every possible event.
\item[Lines 31--32] Return a named tuple and close the function. In \code{up=up}, the name on the left labels the returned field; the value on the right is the quantity just calculated.
\end{description}

\section{Function 3: use a random number to choose the next count}
\lstinputlisting[firstline=35,lastline=44,firstnumber=35]{code/MoranModel.jl}

\begin{description}[style=nextline]
\item[Lines 35--36] Receive the current count, the two change probabilities, and a number \code{u}. Check that the number lies between zero, inclusive, and one, exclusive.
\item[Lines 37--38] If \code{u} is below the increase probability, return a count larger by one.
\item[Lines 39--40] Otherwise test whether \code{u} lies below the combined increase-and-decrease probability. The first branch has already excluded the increase section, so this branch corresponds only to a decrease.
\item[Lines 41--42] All remaining draws mean no change. Return the original count.
\item[Lines 43--44] Close the conditional and the function. Each call returns exactly one integer.
\end{description}

For probabilities \(0.3,0.2,0.5\), draws in \([0,0.3)\) increase the count, draws in \([0.3,0.5)\) decrease it, and draws in \([0.5,1)\) preserve it. Square brackets include an endpoint; round brackets exclude an endpoint. Draws \(0.12\), \(0.41\), and \(0.72\) therefore produce three different outcomes from the \emph{same initial state}.

\begin{workedexample}
\textbf{A sequence needs new probabilities after every event.} The demonstration starts at \(i=2\) and uses the numbers \(0.12,0.72,0.41,0.95,0.10,0.03\). Its successive counts are \(3,3,2,2,3,4\). After the first increase, the population has changed, so the program recalculates its probabilities before using the next draw. Do not apply \(0.3,0.2,0.5\) unchanged to all six events.
\end{workedexample}

\section{Function 4: record an entire population history}
\lstinputlisting[firstline=47,lastline=68,firstnumber=47]{code/MoranModel.jl}

\begin{description}[style=nextline]
\item[Lines 47--49] Receive the random-number generator, population inputs, and the maximum number of events to display. Check both the biological inputs and the nonnegative whole-number horizon. \code{i0} means the initial count; it is a name, not \(i\) multiplied by zero.
\item[Line 50] Allocate storage for integer counts. \code{Vector{Int}} means an array of integers. \code{undef} means the entries are not filled in yet. There are \code{max_events + 1} entries because we also record the initial state.
\item[Lines 51--52] Store the initial count in entry one and copy it into the variable that will track the changing population.
\item[Line 53] Initially no absorption event has been recorded. \code{nothing} is different from zero: zero would mean absorption already occurred at the initial state.
\item[Lines 54--56] If the initial state is already lost or fixed, record absorption at event zero. The double vertical bar inside the \code{if} condition means ``or.''
\item[Line 57] Begin a loop over every event up to the display horizon.
\item[Lines 58--60] If both types are present, calculate this state's probabilities and call \code{next_count} with a fresh draw. The returned count replaces the old value of \code{i}.
\item[Lines 61--63] If this event reaches zero or \code{N}, record the event number. Absorption means the count will remain there in this mutation-free model.
\item[Line 64] Close the condition that permits reproduction updates. If the population was already absorbed, no further random draws were needed.
\item[Line 65] Store the count. After absorption we keep storing the same value, so every displayed path has the same length. This does not change the time of absorption already recorded.
\item[Lines 66--68] Finish the loop and return the array, an absorption flag, and the absorption time. \code{!== nothing} asks whether a time has actually been recorded.
\end{description}

A display horizon is not a rule for deciding fixation. If a path still contains both types at its horizon, it is unfinished. Our separate fixation experiment below has no such horizon.

\section{Function 5: run until loss or fixation}
\lstinputlisting[firstline=71,lastline=81,firstnumber=71]{code/MoranModel.jl}

\begin{description}[style=nextline]
\item[Lines 71--73] Define the function, check the inputs, and initialise the current count.
\item[Line 74] Start an event counter at zero.
\item[Line 75] Continue while the count lies strictly between the two absorbing boundaries.
\item[Lines 76--77] Recalculate the probabilities and use a fresh random draw to update the state.
\item[Line 78] Add one to the number of events. \code{events += 1} is shorthand for \code{events = events + 1}. It also counts events that did not change the allele count.
\item[Line 79] Return to the loop condition. Exit only after reaching a boundary.
\item[Lines 80--81] Return whether the final state is fixation and how many events elapsed. \code{fixed=i == N} labels the field \code{fixed} with a true-or-false answer. There is no cutoff that could mistake a surviving allele for a lost allele.
\end{description}

This function stores only the current state. It does not store all previous states because a fixation experiment needs the final outcome rather than a graph of the whole history. That makes many repetitions use little memory.

\section{Function 6: calculate the exact answer accurately}
\lstinputlisting[firstline=84,lastline=97,firstnumber=84]{code/MoranModel.jl}

Chapter~7 derives the answer. You may read this section once now to see how the code is organised and return after deriving it.

\begin{description}[style=nextline]
\item[Lines 84--85] Define the exact-probability function and check the inputs.
\item[Lines 86--89] A lost allele has fixation probability zero; an already fixed allele has probability one. Return these exact boundary answers immediately.
\item[Lines 90--92] With no selection, return \(i/N\). Handling neutrality separately avoids trying to divide zero by zero in the general expression.
\item[Line 93] Calculate \(\log(1+s)\) using \code{log1p}. The natural logarithm reverses exponentiation by the constant \(e\); Chapter~8 introduces it. The specialised function preserves small differences when \(s\) is tiny.
\item[Line 94] Calculate \(1-\exp[-i\log(1+s)]\). Julia's \code{expm1(y)} computes \(e^y-1\) accurately when \(y\) is close to zero. The leading minus changes its sign, giving \(1-e^y\).
\item[Line 95] Calculate the same expression using \(N\) in place of \(i\).
\item[Lines 96--97] Divide the two quantities and return the exact fixation probability, up to floating-point rounding.
\end{description}

The computer expression is the same formula as
\[
\frac{1-(1+s)^{-i}}{1-(1+s)^{-N}}.
\]
Its appearance differs to protect numerical accuracy. If the computer rounds two very similar numbers to the same stored number, their difference can lose useful digits. This is called cancellation. The specialised functions avoid unnecessary cancellation; they do not add a new biological assumption.

\section{Function 7: check the exact formula by another method}
\lstinputlisting[firstline=100,lastline=121,firstnumber=100]{code/MoranModel.jl}

The backward recurrence links the unknown fixation probability in state \(i\) with those in states \(i-1\) and \(i+1\). Rearranged, it is
\[
(T_i^++T_i^-)u_i-T_i^-u_{i-1}-T_i^+u_{i+1}=0.
\]
Here \(u_i\) is the probability we are trying to find, and \(T_i^+,T_i^-\) are the event probabilities. The boundary answers \(u_0=0,u_N=1\) are known. The remaining \(N-1\) answers must satisfy \(N-1\) linked equations. Chapter~7 solves a two-unknown example by ordinary substitution.

\begin{description}[style=nextline]
\item[Lines 100--101] Define the function and validate population size and selection. We use zero as a valid placeholder count for the shared checker.
\item[Lines 102--105] Create arrays of zeros for the coefficients of the unknowns and the constants on the right side. \code{zeros} stores floating-point zeros. The main diagonal has \(N-1\) entries; each adjacent diagonal has \(N-2\).
\item[Line 106] Visit each nonabsorbing state \(i=1,\ldots,N-1\).
\item[Lines 107--108] Get the event probabilities and store the coefficient of \(u_i\).
\item[Lines 109--111] If a lower interior neighbour exists, store the coefficient of \(u_{i-1}\). For state one the lower neighbour is the known zero boundary, so no unknown entry is required.
\item[Lines 112--116] Store the coefficient of the upper interior neighbour when it exists. At the last interior state, the upper neighbour is \(u_N=1\). Move its term to the right side, giving a positive constant \(T_i^+\).
\item[Lines 117--118] Finish filling the coefficients and assemble a tridiagonal matrix: a compact table with nonzero entries only on the central diagonal and its two neighbours.
\item[Line 119] Julia's backslash asks its linear-algebra tools to solve the linked equations. It is not ordinary division. The vector \code{interior} holds the unknown probabilities in state order.
\item[Lines 120--121] \code{vcat} joins values vertically. Attach the known zero and one boundary values to the interior answers and return all \(N+1\) probabilities.
\end{description}

This calculation does not use the closed-form fixation formula. Agreement checks two different ways of solving the same specified model. It is stronger evidence than running the same formula twice, although it does not prove the model is a realistic description of a particular organism.

\section{Function 8: calculate uncertainty from repeated experiments}
\lstinputlisting[firstline=124,lastline=140,firstnumber=124]{code/MoranModel.jl}

Chapter~6 explains this calculation with numbers before asking you to interpret intervals in the results.

\begin{description}[style=nextline]
\item[Lines 124--126] Accept \code{k} fixations in \code{n} independent repetitions. Require a positive integer number of repetitions and a fixation count from zero through \code{n}.
\item[Line 127] Set the normal-distribution multiplier for a two-sided 95\% interval, approximately \(1.96\).
\item[Line 128] Calculate the observed proportion \(k/n\).
\item[Line 129] Calculate a denominator shared by the next two expressions.
\item[Line 130] Calculate the Wilson interval's centre. \code{2n} is Julia's permitted shorthand for \code{2 * n}.
\item[Line 131] Calculate the distance from that centre to either endpoint. \code{sqrt} means square root. \code{4n^2} means \(4n^2\), not \((4n)^2\).
\item[Lines 132--133] Subtract and add the half-width. \code{max} and \code{min} keep the numerical endpoints within zero and one.
\item[Lines 134--138] Explicitly set the zero-success lower endpoint to zero and the all-success upper endpoint to one. This also removes insignificant floating-point rounding at those boundaries.
\item[Lines 139--140] Return the two named endpoints. An interval is not a guarantee that this particular simulation estimate equals the exact answer.
\end{description}

\section{Function 9: calculate a deterministic comparison}
\lstinputlisting[firstline=143,lastline=155,firstnumber=143]{code/MoranModel.jl}

\begin{description}[style=nextline]
\item[Lines 143--145] Receive population inputs and a horizon and check them.
\item[Line 146] Allocate an array of floating-point frequencies, with room for the initial value.
\item[Lines 147--148] Convert the initial count into a frequency and store it.
\item[Line 149] Step through the same number of events as the stochastic display.
\item[Line 150] Calculate the expected one-event frequency change evaluated at the current deterministic frequency. Multiplying out the expression gives \(sx(1-x)/[N(1+sx)]\).
\item[Line 151] Add this change to the current frequency. There is no random draw.
\item[Lines 152--155] Store the updated frequency, continue the loop, and return the array.
\end{description}

At \(N=5,i=2,s=0.5\), the first change is \((0.3-0.2)/5=0.02\). The deterministic frequency changes from \(0.40\) to \(0.42\). A real model event can only change it to \(0.60\), \(0.20\), or leave it at \(0.40\). The deterministic number is an expectation-based comparison, not a possible outcome of that one event. Chapter~8 explains why repeatedly applying it is a mean-field approximation.

\begin{checkpoint}
\textbf{Try it.} Why do we increment the event counter after a no-change event? Why can the fixation function avoid keeping an array?\\[3pt]
\textbf{Answer.} Reproduction and replacement still occurred, so one model event elapsed. The count happened to remain the same. For a fixation experiment we need only whether the allele eventually fixed and how long the experiment took; earlier counts are unnecessary unless we want to draw a history.
\end{checkpoint}
